%!TEX root = ../template.tex
%% ====================================================================
%% Capitulo 1 -- Introducao
%% v1 2026-09-13. v2 2026-09-17: adaptado ao levantamento de requisitos
%% de fevereiro de 2026 (random/Perguntas Contextualizacao Problema.docx),
%% ao titulo "Medir nao basta" e aos comentarios C1-02..C1-07. Versao
%% anterior em _revisao/2026-09-17-cap1-antes.tex.
%%
%% Regras desta versao: AO90; impessoal com -se; sem numeros de
%% resultados (so contagens estruturais); sem mencao ao artigo; nenhum
%% nome de pessoa; nada do que o levantamento diz "nao apurado" vira
%% afirmacao (politica energetica, obrigacoes legais, nao-conformidades,
%% tempo dos relatorios). Normalizacao: "so pela carga", nunca "nao ha".
%% Marcadores: %% FIGURA (sugestao de figura/esquema).
%%
%% Setup narrativo (skill narrativa, modo nao-ficcao):
%%   focalizacao -- o problema e de quem opera o SGE (a equipa de
%%     monitorizacao, o tecnologo, o auditor), nunca do autor;
%%   altitude -- banda d2-d3, uma descida por seccao: 1.1 a rotina da
%%     equipa e os permutadores; 1.2 o esforco antes da auditoria; 1.5 a
%%     fronteira da unidade;
%%   informacao -- conclusao a frente ("medir nao basta" dito na primeira
%%     pagina); tensao pelas apostas: fraquezas nomeadas antes do revisor
%%     (SQ4 retirada, extracao sem verificacao humana, criterios de sucesso
%%     nunca definidos, dados ja explorados);
%%   recorrencia -- a pergunta de abertura (consome hoje mais do que a
%%     referencia preve para a carga que tem?) volta no fecho de 1.6 com a
%%     resposta a trazer o seu proprio estado; "mede-se tudo, sabe-se
%%     tarde" (1.1) volta como "separa medir de saber" (1.6); "duas
%%     metades de um requisito" plantado em 1.1, volta em 1.4 sem ser
%%     nomeado; os permutadores em 1.1 voltam em 1.5 como razao da CDU.
%% ====================================================================

\chapter{Introdução}
\label{cha:introducao}
% [REV-2026-09-13][C1-01][P2][ESTRUTURA]
\revisaotese{C1-01}{P2}{ESTRUTURA}
{Estado da dissertação e ordem de revisão}
{Revisão crítica de 2026-09-13, atualizada a 2026-09-17. O arco plataforma
-\textgreater{} diagnóstico -\textgreater{} proposta é defensável. Faltam ainda: as duas
secções do Corpus~B e o \emph{gap} no Cap.~3; o Cap.~8; os resumos; os
Apêndices B--D; a redução do Apêndice~E ao que o corpo cita. Não
apresentar como entregues elementos ainda em esqueleto. Procurar
REV-2026-09-13 em todos os .tex.}%
% [/REV-2026-09-13]

\section{Contexto e motivação}
\label{sec:contexto}

A Refinaria de Sines mede a sua energia com um detalhe que poucas
organizações têm. Cada caudal de combustível e de vapor, cada temperatura
e cada poder calorífico ficam registados no historiador de processo, ao
minuto; os balanços mensais de produção e de utilidades fecham-se e
arquivam-se; a eletricidade é medida, faturada e repartida pelas
unidades. A refinaria é certificada segundo a ISO~50001, tem um sistema
de gestão de energia implementado e quatro indicadores de desempenho
energético definidos: consumos específicos por unidade, índice de
intensidade energética, emissões de CO$_2$ associadas ao consumo, e as
retas de orçamento que servem de referência ao consumo de cada
unidade.\footnote{As afirmações deste capítulo sobre a rotina anterior à
plataforma provêm de um levantamento de requisitos junto da orientação
industrial, feito em fevereiro de 2026, no arranque do estágio: trinta e
seis perguntas em cinco blocos (contexto organizacional, estado atual,
lacunas, restrições técnicas, expectativas), com respostas registadas por
escrito. Cita-se como \enquote{levantamento} ao longo do texto. O que
nele ficou por apurar não é usado.} E, contudo, quando se pergunta se uma
unidade consome hoje mais do que a referência prevê para a carga que tem,
a resposta não sai de nenhum destes sistemas.

Sai da equipa de monitorização de processo, e sai devagar. O
acompanhamento dos usos significativos de energia faz-se em folhas de
registo que o levantamento descreve como pesadas e pouco acessíveis; a
análise faz-se em Excel, à mão, sem ferramenta dedicada nem
\emph{dashboard}; o resultado chega à gestão num relatório trimestral,
também manual. Quando um desvio é detetado, investiga-o o tecnólogo da
unidade, sem procedimento formal que o levantamento tenha apurado. Os
problemas de qualidade de dados são regulares e detetam-se com atraso,
porque nada os valida automaticamente. Os dados existem em três lugares
que não se falam --- o historiador, os livros de balanço e os ficheiros
de cálculo das equipas --- e a monitorização vive na fronteira entre
eles, reconciliada caso a caso. O próprio levantamento descreve a gestão
de energia como orientada para o cumprimento dos requisitos de auditoria,
e não para a melhoria contínua. Mede-se tudo; sabe-se tarde.

O pedido que deu origem a este trabalho foi, na sua formulação inicial,
um \emph{dashboard}: uma plataforma que juntasse as fontes, calculasse
os indicadores que a equipa já produzia e os mostrasse por unidade, por
vetor energético e por dia. É um pedido legítimo, e foi cumprido. Resolve
a primeira metade do problema: a informação passa a existir num modelo
único, validado contra os números oficiais, atualizado todas as noites.
Mas medir melhor, e mais depressa, ainda não é saber. A pergunta que o
\emph{dashboard} esconde é outra, e é a desta dissertação: \emph{o que é
preciso para que uma plataforma de monitorização energética seja
fiável?}

Fiável não quer dizer infalível. Nenhum detetor estatístico garante que
cada alarme corresponde a um problema e que a ausência de alarme
corresponde à ausência de problema, e esta dissertação não mede tal
coisa. Fiável quer dizer quatro coisas verificáveis. Cada número é
rastreável até à fonte e à regra que o produziu. O domínio em que a
comparação é válida está declarado, e sabe-se quando se está fora dele. O
desempenho da referência está quantificado: contra que comparador, com
que erro, em que período. E a incerteza acompanha o resultado, em vez de
ficar implícita numa reta desenhada sobre uma nuvem de pontos. A
definição é exigente por ser modesta: não promete alarmes certos, promete
que se sabe o que cada alarme vale.

A fiabilidade assim definida tem duas camadas, e a segunda é a que
costuma faltar. A primeira é a dos dados: a plataforma tem de dizer a
verdade sobre o que a unidade consumiu e processou, o que exige as fontes
reconciliadas num modelo único, com cada divergência explicada. A segunda
é a da analítica: a plataforma compara essa verdade com uma referência
--- uma \emph{baseline} energética --- que diz quanto a unidade
consumiria, àquela carga, em condições de referência. Não é o consumo
ótimo, nem o que a unidade \enquote{devia} consumir: uma regressão sobre
o histórico estima o que foi habitual, não o que é termodinamicamente
possível. Se a referência estiver errada, os dados podem estar certos e o
alarme continuar a não significar nada.

As duas camadas não são uma escolha de desenho; são as duas metades de
um requisito. A ISO~50001 obriga a organização a monitorizar, medir e
analisar o seu desempenho energético e a responder a desvios
significativos~\cite[\S9.1]{ISO50001_2018}; a ISO~50006 operacionaliza a
comparação através do par indicador--referência, os indicadores de
desempenho energético (EnPI) e os consumos energéticos de referência
(EnB, as \emph{baselines})~\cite{ISO50006_2014}. A infraestrutura de
medição é a primeira metade; a \emph{baseline} validada é a segunda. Uma
organização pode ter a primeira e não ter a segunda, e é o que se
encontrou em Sines: um método de \emph{baseline} aplicado com disciplina
há seis anos, reproduzível ao dígito, que normaliza o consumo por uma
única variável, a carga --- nenhuma variável externa, nenhum fator
estático entra --- e cuja adequação ao fim que a norma lhe atribui nunca
tinha sido avaliada com um protocolo declarado.

% [REV-2026-09-13][F01][P2][FIGURA]
\revisaotese{F01}{P2}{FIGURA}
{Priorizar o esquema das duas camadas}
{Aproveitar esta sugestão existente como única figura da Introdução.
Local: depois do parágrafo que apresenta as duas camadas. Desenhar fontes
-\textgreater{} reconciliação -\textgreater{} indicador -\textgreater{} referência -\textgreater{} desvio -\textgreater{} investigação,
com duas faixas (dados/analítica) e os capítulos correspondentes.
Distinguir graficamente entrega em produção de proposta ainda por testar.
Meia página; desenho vetorial próprio; legenda explica que rastreabilidade
não equivale a exatidão física nem um desvio identifica a sua causa.
Dispensar a Figura 1.3 se repetir este mapa.}%
% [/REV-2026-09-13]
%% FIGURA 1.1 (sugestao): esquema das duas camadas da fiabilidade.
%% Caixa superior "camada de dados" (historiador + balancos + ficheiros
%% -> modelo dimensional unico) ligada a ISO 50001 §9.1 e ao Cap. 5;
%% caixa inferior "camada analitica" (baseline -> desvio -> alarme)
%% ligada a ISO 50006 (EnPI/EnB) e ao Cap. 6. Seta de retorno do Cap. 7
%% para a camada analitica ("bloco de saude"). Desenho proprio, sem
%% dados; serve de mapa mental para o resto do capitulo.

Isto importa por uma razão que a própria orientação industrial
identificou antes de existir uma linha de código. Na unidade de
destilação, os ciclos de limpeza dos permutadores da bateria de
pré-aquecimento alteram de forma previsível o consumo específico do forno
a jusante: a incrustação acumula-se, o forno compensa, a limpeza repõe
parte do que se perdeu, e o ciclo recomeça. O estado de incrustação é uma
variável relevante no sentido da ISO~50006 --- afeta o consumo e muda
rotineiramente --- mas não é medido, e a referência não o vê. O
levantamento reconheceu a tensão: as métricas precisam de uma referência
robusta, e uma referência rígida não serve todas as situações, pelo que
se adapta a análise \enquote{caso a caso}, com \enquote{espírito
crítico}. Uma referência que se reajusta todos os anos resolve a mesma
tensão de outra maneira, pior: em cada janeiro absorve a deriva do ano
anterior e volta a zero. Sobram para os alarmes os desvios rápidos; os
lentos, que são os que a gestão de energia existe para apanhar, ficam
invisíveis por construção. Medir não basta. Esta dissertação parte do
pedido de um \emph{dashboard} e chega a essa conclusão pelo caminho mais
longo, o de mostrar o que falta.

\section{Formulação do problema}
\label{sec:problema}

O levantamento descreve o problema sem rodeios, e é por aí que se
começa. Não há acompanhamento regular do desempenho energético: o ideal
seria diário, e o que existe concentra-se nos períodos que antecedem as
auditorias. Há decisões de melhoria que se tomam com atraso, por falta de
visibilidade atempada sobre os dados. A refinaria consegue demonstrar
melhoria com dados, mas o processo é descrito como pouco eficiente e
lento; quanto tempo custa cada relatório trimestral nunca foi
quantificado, e não se inventa aqui um número. Nenhuma variável externa
entra na normalização dos indicadores. Em fevereiro de 2026, antes de
qualquer resultado, a leitura da norma contra o levantamento já apontava
essa ausência como lacuna face aos requisitos de determinação de
variáveis relevantes e de normalização~\cite[\S6.4--6.5]{ISO50001_2018}.
Esta parte do problema, uma plataforma resolve-a em grande medida, e o
Capítulo~\ref{cha:plataforma} mostra como.

Por trás dela há duas razões, e a segunda explica por que motivo a
primeira não se resolve com mais dados.

A primeira é uma assimetria normativa. A ISO~50001 é uma norma de
requisitos: diz o que deve existir (indicadores, \emph{baselines},
normalização quando há variáveis relevantes, resposta a desvios) e é
certificável sobre isso. A ISO~50006 e a ISO~50015 são normas de
orientação: descrevem como construir uma \emph{baseline} e como medir e
verificar, sem carácter vinculativo~\cite{ISO50006_2014,ISO50015_2014}.
Nenhuma das três fixa um critério pelo qual uma \emph{baseline} é aceite
ou rejeitada. A organização é obrigada a comparar o desempenho com a
referência e a reagir ao que vê; não é obrigada a demonstrar que a
referência merece a comparação. O Capítulo~\ref{cha:enquadramento}
percorre esta assimetria cláusula a cláusula e mostra que o silêncio é
uma escolha de desenho normativo, não uma impossibilidade técnica:
protocolos de medição e verificação nascidos no setor dos edifícios
fixam critérios quantitativos, e a família ISO cita-os sem os adotar.

A segunda é um pressuposto que as normas fazem sem o enunciar.
Normalizar o consumo pela carga, comparar períodos \enquote{em condições
equivalentes} e alarmar sobre um desvio só faz sentido se a relação
entre consumo e variáveis relevantes for estável no tempo --- se o que
a unidade consumia a uma dada carga no período de referência for o que
consome a essa carga hoje. É um pressuposto de estacionariedade da
relação, não da série, e é razoável em muitos contextos: um edifício,
uma linha de montagem. Em processo contínuo não pode ser assumido; tem
de ser verificado. A relação consumo--carga move-se por razões que
convém distinguir, porque pedem respostas diferentes: a carga muda de
composição com cada crude; o equipamento degrada-se, como a bateria de
pré-aquecimento que incrusta e obriga o forno a compensar; a medição
altera-se quando um instrumento é recalibrado ou substituído; e as
intervenções --- uma limpeza, uma paragem --- repõem parte do que se
perdeu. Nem tudo isto é deriva estrutural e nem tudo é ineficiência: a
incrustação é, ao mesmo tempo, um mecanismo físico previsível e uma
perda de eficiência que a operação pode antecipar ou adiar. Um modelo
estático, ajustado sobre um período fixo, não distingue nenhuma destas
coisas: vê apenas resíduos. E se o modelo for reajustado todos os anos,
deixa de ver sequer os resíduos.

%% FIGURA 1.2 (sugestao, opcional): esboco conceptual, sem dados, de
%% consumo vs carga com a relacao a deslocar-se ano a ano e o refit anual
%% a "apanhar" cada deslocacao (tres retas paralelas, cada uma com a sua
%% nuvem; a banda de alarme de cada ano centrada na sua propria reta).
%% Legenda: "o que a recalibracao anual absorve". Se for usada, tem de
%% ser explicitamente conceptual, para nao antecipar a Figura real do
%% Cap. 6 (sec:estabilidade). Alternativa: deixar so para o Cap. 2 §2.4.

O que está em jogo é operacional e normativo ao mesmo tempo.
Operacional, porque uma \emph{baseline} que não distingue deriva de
ineficiência produz alarmes que a operação aprende a ignorar e
silêncios que a gestão aprende a ler como normalidade. Normativo, porque
uma organização pode estar em conformidade formal com a ISO~50001 (tem
EnPI, tem EnB, normaliza, reporta) e não ter forma de saber se as
suas \emph{baselines} cumprem alguma das três funções que lhes são
pedidas: normalizar o consumo pela carga, sustentar comparações entre
períodos e alimentar alarmes. A certificação não pergunta; a norma não
exige que se pergunte.

% [REV-2026-09-13][C1-04][P1][ARGUMENTO]
\revisaotese{C1-04}{P1}{ARGUMENTO}
{Fechar a lacuna com a síntese do Cap.~3}
{O parágrafo seguinte está na formulação provisória. Depois da síntese do
Corpus~B: delimitar a lacuna por âmbito, período, fontes pesquisadas e
características que as abordagens existentes cobrem; confirmar contra a
síntese final a frase sobre a quase ausência da incrustação na literatura
(leitura assistida, sem verificação humana por ocorrência). A inovação
pode ser a integração e demonstração industrial do procedimento, mesmo
que os componentes já existam.}%
% [/REV-2026-09-13]
É este o problema, na formulação que o Capítulo~\ref{cha:revisao}
delimita: \emph{falta demonstrar, no caso estudado, que as
\emph{baselines} em produção são adequadas aos usos que a norma lhes
atribui, e falta um procedimento operacional para o fazer; sem ele, uma
organização não tem meio de saber se a sua plataforma de monitorização é
fiável na segunda camada, por mais fiável que seja na primeira.} A
revisão sistemática do âmbito pergunta à literatura o que existe.
Adianta-se aqui só o que muda a maneira de ler o resto do documento: os
estudos em processo contínuo reportam sobretudo sintomas no modelo ---
relações que não se mantêm, indicadores que deixam de ter significado,
grandezas que não são medidas --- e quase nunca a causa física por trás
deles. A incrustação que a orientação industrial identificou de imediato
é, nessa literatura, quase ausente. Os capítulos empíricos mostram o que
o problema custa numa refinaria real e o que é preciso para o tratar.

\section{Objetivos e questões de investigação}
\label{sec:objetivos}

A pergunta central já foi enunciada: o que é preciso para que uma
plataforma de monitorização energética seja fiável, nas duas camadas.
Responde-se-lhe com três objetivos encadeados, em que cada um é a
condição do seguinte.

O primeiro é \emph{construir}: reunir as três fontes de dados da
refinaria num modelo dimensional único; validar que esse modelo
reproduz os números oficiais, célula a célula, com cada divergência
explicada; calcular sobre ele as métricas energéticas que a equipa de
monitorização usa; e entregar o resultado em produção. Este objetivo
responde ao pedido industrial e produz o instrumento.

O segundo é \emph{diagnosticar}: replicar exatamente, dentro desse
instrumento, o método de \emph{baseline} que a refinaria usa há seis
anos, e avaliá-lo com um protocolo estatístico especificado e selado
antes de qualquer resultado: poder preditivo contra referências
triviais, estabilidade dos parâmetros, calibração das bandas de alarme,
detetabilidade e suficiência da variável explicativa, com correção de
multiplicidade e análise de sensibilidade às convenções herdadas do
método. A replicação exata é o que garante que se critica o método real
e não uma reconstrução dele.

O terceiro é \emph{prescrever}: converter o diagnóstico numa proposta,
linhas orientadoras para construir, aceitar e manter \emph{baselines}
em processo contínuo, com um critério de aceitação declarado antes dos
resultados; demonstrá-la na própria unidade; e devolvê-la à plataforma
como especificação do que o \emph{dashboard} tem de mostrar para que uma
\emph{baseline} seja usada com conhecimento do seu estado.

Os três objetivos são acompanhados por uma revisão sistemática do
âmbito da literatura, registada e reportada segundo a
PRISMA-ScR~\cite{Tricco2018}, cuja pergunta primária é em que medida os
modelos de medição do desempenho energético propostos na literatura para
cumprir o enquadramento ISO~50001/50006 são adequados à indústria de
processo contínuo, e que adaptações metodológicas são necessárias para a
sua aplicação. Três subquestões a decompõem, e são também subquestões
desta dissertação: a literatura responde-lhes a partir do que está
publicado; os capítulos empíricos, a partir de um caso real.

\begin{description}
  \item[SQ1] Que modelos de medição do desempenho energético existem no
    âmbito do enquadramento ISO~50001/50006, para que contextos foram
    desenvolvidos e validados, e como se distribuem por tipo de
    organização? --- respondida pela análise bibliométrica do Corpus~A
    (Capítulo~\ref{cha:revisao}).
  \item[SQ2] Que mecanismos de variabilidade operacional, específicos da
    indústria de processo contínuo, podem tornar esses modelos
    inadequados para uma monitorização contínua alinhada com as normas
    de referência? --- respondida pela síntese narrativa do Corpus~B e,
    empiricamente, pelo diagnóstico do Capítulo~\ref{cha:diagnostico}: o
    caso confirma ou contradiz o que a literatura antecipa.
  \item[SQ3] Como podem os modelos de EnPI e EnB ser adaptados à
    variabilidade operacional dos processos contínuos, dentro dos
    requisitos das normas de referência? --- respondida pelas
    adaptações identificadas no Corpus~B e pelas linhas orientadoras do
    Capítulo~\ref{cha:discussao}.
\end{description}

O protocolo registava uma quarta subquestão, sobre a aplicabilidade das
recomendações identificadas numa refinaria real. Foi retirada a 16 de
setembro de 2026, porque a verificação humana da extração que a
sustentaria não foi executada; a retirada está declarada no
Apêndice~\ref{app:protocolo-sr} (desvio D8). A aplicabilidade é tratada
apenas pelo caso empírico, e a revisão não faz afirmações de
transferibilidade para a refinação. A Tabela~\ref{tab:objetivos-perguntas}
liga cada objetivo às perguntas que o concretizam, à evidência que as
responde e ao estado em que essa resposta se encontra.

\begin{table}[htbp]
  \centering\small
  \caption{Objetivos, perguntas empíricas, evidência e estado da resposta.
  As perguntas da revisão constam para completar o mapa; a revisão não
  responde a nenhuma pergunta empírica.}
  \label{tab:objetivos-perguntas}
  \begin{tabularx}{\linewidth}{@{}>{\raggedright\arraybackslash}p{1.7cm}>{\raggedright\arraybackslash}X>{\raggedright\arraybackslash}X>{\raggedright\arraybackslash}p{3.1cm}@{}}
    \toprule
    Objetivo & Pergunta & Evidência e onde & Estado \\
    \midrule
    Construir & Os dados reconciliados reproduzem os números oficiais, e cada divergência tem explicação? & Validação célula a célula contra o fecho mensal (Secção~\ref{sec:validacao-dados}) & respondida \\
    \addlinespace
    Diagnosticar & A plataforma reproduz ao dígito o método de \emph{baseline} em produção? & Replicação exata sobre o modelo de dados (Secção~\ref{sec:replicacao}) & respondida \\
    Diagnosticar & A referência prevê melhor do que comparadores triviais, por vetor e por ano? & Contraste de perdas emparelhado com intervalos por blocos (Secção~\ref{sec:poder-preditivo}) & respondida nos quatro vetores diários; a eletricidade fica sem diagnóstico diário \\
    Diagnosticar & A relação consumo--carga é estável entre anos? & Contrastes entre anos com covariância robusta à dependência (Secção~\ref{sec:estabilidade}) & respondida \\
    Diagnosticar & A banda de alarme cobre o que promete, no conjunto e por estrato de carga? & Cobertura marginal e por estrato (Secção~\ref{sec:calibracao}) & respondida \\
    Diagnosticar & Que desvio o método consegue detetar, e a carga chega como variável explicativa? & Detetabilidade e modelos encaixados com covariáveis de processo (Secções~\ref{sec:detetabilidade} e~\ref{sec:suficiencia}) & respondida \\
    \addlinespace
    Prescrever & Um candidato construído pelas linhas orientadoras passa o critério de aceitação fixado antes dos resultados? & Demonstração nos mesmos pares anuais (Secção~\ref{sec:demonstracao}) & respondida; retrospetiva \\
    Prescrever & O que deteta um monitor sequencial, e com que atraso, quando a verdade é conhecida? & Simulação sobre referência congelada (Secção~\ref{sec:verdade-conhecida}) & respondida por simulação \\
    Prescrever & O que tem o \emph{dashboard} de mostrar para que uma \emph{baseline} seja usada com conhecimento do seu estado? & Especificação do bloco de saúde (Secção~\ref{sec:implicacoes-plataforma}) & proposta; não implantada \\
    \addlinespace
    Revisão, SQ1 & Como se distribuem os modelos por setor, organização e enquadramento normativo? & Bibliometria do Corpus~A (Secção~\ref{sec:corpus-a}) & respondida \\
    Revisão, SQ2--SQ3 & Que mecanismos tornam os modelos inadequados, e que adaptações se propõem? & Síntese narrativa do Corpus~B (Secção~\ref{sec:corpus-b}) & em redação \\
    Revisão, SQ4 & As recomendações são aplicáveis numa refinaria? & --- & retirada (Apêndice~\ref{app:protocolo-sr}, D8) \\
    \bottomrule
  \end{tabularx}
\end{table}
%% Tabela 1.1: atualizar a coluna "Estado" (SQ2--SQ3) quando o Cap. 3
%% fechar. Nao apresentar como respondido o que ainda esta em esqueleto.

O desenho de investigação assenta em três rotas de evidência, cada uma
com a sua regra fixada antes do resultado e com a cronologia declarada,
para que se saiba o que se sabia em cada momento. A revisão diz o que
existe: o protocolo foi registado publicamente a 24 de fevereiro de 2026,
depois das pesquisas nas bases bibliográficas e antes da recuperação dos
textos e de qualquer extração; cada alteração posterior é uma emenda
datada, e os desvios que não foram objeto de emenda estão declarados com
a mesma visibilidade (Apêndice~\ref{app:protocolo-sr}). O caso real diz o
que acontece: o método de produção é replicado com exatidão, sobre os
dados da própria refinaria, antes de ser avaliado. O protocolo de
avaliação diz como se sabe: os estimandos, os testes, os limiares e a
correção de multiplicidade foram fixados e versionados antes de se olhar
para um resultado, na lógica do pré-registo~\cite{Nosek2018}, e
submetidos a quatro revisões adversariais produzidas por modelos de
linguagem e adjudicadas pelo autor --- um controlo de defeitos, não uma
revisão por pares. Os dados, esses, não eram desconhecidos: a exploração
feita ao construir a plataforma antecedeu o protocolo, e o
Capítulo~\ref{cha:metodos} diz onde isso pesou. Uma grelha de
sensibilidade reavalia as conclusões sob variações das convenções que o
método herdou.

\section{Contribuições}
\label{sec:contribuicoes}

As contribuições dividem-se pelas duas metades do requisito. As
primeiras pertencem à refinaria. As segundas transferem-se para quem
opere \emph{baselines} lineares anuais sobre a carga, com série diária
de pelo menos um ano: é a classe em que foram demonstradas, e a única
sobre a qual se afirma alguma coisa.

Para a refinaria, entregaram-se duas coisas e especificou-se uma
terceira. Uma plataforma de monitorização energética em produção, sobre
um modelo dimensional único que reconcilia as três fontes e cobre
dezasseis unidades processuais em três fábricas, validada contra os
números oficiais célula a célula (Capítulo~\ref{cha:plataforma}). Um
\emph{dashboard} desenhado como percurso de diagnóstico, da refinaria
para a unidade, da unidade para o vetor, do vetor para a \emph{baseline}
que o julga (Secção~\ref{sec:dashboard}). E a especificação de um bloco
de saúde das \emph{baselines}: duas pistas em vez de uma, domínio de
validade visível, calibração monitorizada em vez de assumida, uma ficha
por \emph{baseline} e políticas de estado registadas, formulada como
requisitos, cada um ancorado num resultado, e ainda por implantar
(Secção~\ref{sec:implicacoes-plataforma}).

As contribuições metodológicas são quatro. Um protocolo de avaliação de
\emph{baselines} energéticas, especificado antes dos resultados, com
correção de multiplicidade e grelha de sensibilidade, reutilizável na
classe acima (Secção~\ref{sec:protocolo-baselines}). O diagnóstico do
método em produção: cinco modos de falha qualitativamente distintos, um
por vetor energético, e o mecanismo que os liga --- a recalibração anual
absorve a deriva que devia expor, e o sistema fica cego por construção
aos fenómenos lentos e ruidoso nos rápidos
(Capítulo~\ref{cha:diagnostico}). Linhas orientadoras para
\emph{baselines} em processo contínuo, com um critério de aceitação
selado e demonstradas na unidade, incluindo o que muda quando se congela
uma referência em vez de a reajustar e o que um monitor sequencial
consegue detetar quando a verdade é conhecida por simulação
(Capítulo~\ref{cha:discussao}). E uma revisão sistemática do âmbito
registada, com corpus fechado e extração assistida por modelos de
linguagem sobre um instrumento congelado, sem verificação humana por
ocorrência, que situa o caso na literatura no estado em que o
Capítulo~\ref{cha:revisao} a documenta (Apêndices~\ref{app:protocolo-sr}
e~\ref{app:validacao}).

As quatro últimas não existiriam sem as três primeiras, e as três
primeiras ficariam por acabar sem as quatro últimas: o instrumento foi
construído para tornar o diagnóstico possível, e o diagnóstico
devolve-lhe a especificação de que ele precisava.

\section{Âmbito e enquadramento do trabalho}
\label{sec:ambito}

O trabalho decorreu num estágio curricular na Refinaria de Sines da
Galp, entre fevereiro e julho de 2026, integrado na equipa de
monitorização de processo, com orientação académica na NOVA FCT e
orientação industrial na refinaria. Não teve financiamento específico:
os custos de computação da revisão e do diagnóstico foram suportados
pelo autor, e o estágio é o único vínculo do autor à organização. A
plataforma foi construída, validada e entregue nesse período; o
diagnóstico e a proposta foram concluídos depois, sobre uma extração
completa e verificada dos dados e do modelo, feita antes da saída. A
dissertação é objeto de um pedido de confidencialidade, com embargo até
setembro de 2029. Os nomes dos sistemas de informação, as designações
das unidades, os coeficientes e os valores absolutos mantêm-se, porque
são eles que tornam o trabalho verificável; as pessoas são identificadas
pelo cargo.

As fronteiras são quatro, todas escolhidas; a primeira foi escolhida em
conjunto com a orientação industrial.

\emph{Uma unidade.} O caso é a unidade de destilação atmosférica (CDU;
na designação interna, CC), com os seus cinco vetores energéticos (fuel
gás, três níveis de vapor e eletricidade), quatro dos quais com
resolução diária; a eletricidade é uma alocação mensal e fica fora de
qualquer inferência diária. O âmbito foi acordado no arranque: uma
unidade primeiro, expansão às restantes depois de validado o piloto, que
não foi desenhado como prova de conceito descartável mas como base de
uma ferramenta a usar em produção. A expansão aconteceu na camada de
dados --- a plataforma cobre dezasseis unidades --- e não no diagnóstico,
que ficou na CDU. A escolha tem três razões. A unidade processa todo o
crude e é, pela sua dimensão, um uso significativo de energia por si
só. Foi para ela que a orientação industrial deu o exemplo dos
permutadores: se um método de \emph{baseline} estático falha, é aqui que
falha de forma legível. E a resolução acordada --- atualização diária,
sem requisito de tempo real, com pelo menos um ano completo de histórico
para qualquer referência --- é a resolução a que quatro dos seus cinco
vetores estão disponíveis.

\emph{A classe linear estática como objeto.} O método de produção é uma
regressão linear anual do consumo sobre a carga, e é essa classe que se
diagnostica e que as linhas orientadoras levam ao máximo defensável:
domínio explícito, exclusões por causa, estimação robusta, bandas que
respeitam a dependência, critério de aceitação. Modelos com estado de
degradação, que seriam a resposta completa à não-estacionariedade,
ficam fora: o Capítulo~\ref{cha:conclusoes} traça o roteiro para eles e
explica por que razão os critérios de aceitação aqui construídos são a
condição para lá chegar.

\emph{Sem inferência causal.} Nenhum desvio é classificado como
ineficiência e nenhum alarme como falso, porque isso exige o calendário
operacional, os rótulos de eventos e as margens materiais elicitadas
com a operação, que não estavam disponíveis. A dissertação mede o que o
método vê e o que não vê; não diz o que a unidade devia ter feito.

\emph{Monitorização contínua, não medição e verificação de projetos.}
Os protocolos de M\&V tratam de provar as poupanças de uma intervenção;
aqui o objeto é a \emph{baseline} que julga cada dia de operação, um
problema distinto que partilha instrumentos com o primeiro.

Estas fronteiras têm um custo, e declara-se já. Uma unidade e um único
\emph{snapshot} de dados já explorados não permitem afirmações fora de
amostra, e a demonstração das linhas orientadoras é retrospetiva. A
eletricidade fica sem diagnóstico diário. O que a recalibração anual
mascarou em cada ano não é mensurável sem o calendário operacional. Os
critérios de sucesso da ferramenta, que o levantamento deixou por
definir, nunca chegaram a ser fixados: a plataforma foi avaliada pela
reconciliação e pelo diagnóstico que tornou possível, não contra
métricas acordadas com quem a usa. A articulação com o sistema integrado
de gestão da refinaria ficou fora de âmbito por decisão conjunta. A
Secção~\ref{sec:limitacoes} desenvolve cada uma destas limitações;
enunciam-se aqui para que se saiba desde o início o que este trabalho
não pode dizer.

\section{Estrutura do documento}
\label{sec:estrutura}
% [REV-2026-09-13][C1-08][P3][REDACAO]
\revisaotese{C1-08}{P3}{REDACAO}
{Reduzir repetição entre abertura, contribuições e roteiro}
{A Introdução repete várias vezes o reajuste que absorve a deriva, a
fiabilidade em duas camadas e o percurso dos capítulos. Objetivo editorial:
reduzir cerca de um quarto depois de estabilizar as perguntas, conservando
problema, lacuna delimitada, objetivos, âmbito e contribuições.
Deixar a demonstração do mecanismo para os Caps. 6/7.}%
% [/REV-2026-09-13]

A dissertação lê-se em três andamentos, instrumento, diagnóstico e
prescrição, precedidos por dois capítulos que estabelecem o que as
normas exigem e o que a literatura oferece.

O Capítulo~\ref{cha:enquadramento} reconstrói a cadeia de medição
implícita na família ISO~50001 e mostra, cláusula a cláusula, onde a
norma exige, onde deixa à organização e que pressuposto tácito cada elo
importa para funcionar; introduz a especificidade dos processos
contínuos e apresenta o caso de estudo. O Capítulo~\ref{cha:revisao}
pergunta à literatura se o problema está resolvido: protocolo e
registo, estratégia de pesquisa e seleção, análise bibliométrica do
Corpus~A, síntese narrativa do Corpus~B e a formulação do \emph{gap}
que o resto do documento trata.

O Capítulo~\ref{cha:metodos} descreve como se construiu e como se
avaliou: fontes e ingestão, modelo dimensional e reconciliação, cálculo
de métricas, o protocolo de avaliação estatística das \emph{baselines}
e a cadeia de proveniência que torna cada número reproduzível. O
Capítulo~\ref{cha:plataforma} é o instrumento: a arquitetura entregue,
a prova de que a camada de dados diz a verdade, o \emph{dashboard} e o
que ficou de fora. O Capítulo~\ref{cha:diagnostico} executa-se sobre o
modelo de dados do capítulo anterior (é, literalmente, uma consulta a
ele) e audita a camada analítica: replicação exata, perfil dos dados e
domínio de validade, as quatro dimensões do protocolo e os diagnósticos
complementares, fechados numa síntese dos modos de falha. O
Capítulo~\ref{cha:discussao} lê o diagnóstico em conjunto, tira as
implicações para a conformidade normativa, propõe as linhas
orientadoras, demonstra-as na unidade, devolve-as à plataforma como
especificação e declara os limites do que afirma. O
Capítulo~\ref{cha:conclusoes} conclui e traça o roteiro dos modelos
sucessores.

%% FIGURA 1.3 (sugestao): mapa de leitura. Tres colunas (instrumento /
%% diagnostico / prescricao) com os Caps. 5, 6 e 7; por cima, Caps. 2 e
%% 3 como "o que as normas exigem" e "o que a literatura oferece"; por
%% baixo, Cap. 4 como base comum; setas: 5->6 ("consulta ao modelo de
%% dados"), 6->7 ("diagnostico -> prescricao"), 7->5 ("bloco de saude").
%% Apendices como faixa lateral. Uma figura, meia pagina, sem texto
%% corrido dentro das caixas. Dispensar se a Figura 1.1 a tornar redundante.

Os apêndices guardam o que o corpo cita e não reproduz. Da revisão, o
protocolo e as emendas (Apêndice~\ref{app:protocolo-sr}) e a extração
assistida e a sua validação (Apêndice~\ref{app:validacao}). Do diagnóstico,
as tabelas completas (Apêndice~\ref{app:tabelas}). Da plataforma, o manual
do \emph{dashboard}, que inclui o dicionário de dados e o mapeamento mestre
(Apêndice~\ref{app:manual}). Por fim, as definições matemáticas dos
instrumentos de avaliação (Apêndice~\ref{app:matematica}).

A pergunta com que o capítulo abriu --- consome a unidade hoje mais do
que a referência prevê para a carga que tem? --- volta no fim do
documento com uma resposta diferente da que o \emph{dashboard} dá: um
número acompanhado do seu estado --- em que domínio a referência é
válida, contra que comparador foi medida, com que cobertura, desde
quando está em vigor e por que razão ainda não foi substituída. É isso que o bloco de saúde passaria a trazer, e é isso que
separa medir de saber.
